\documentclass[10pt]{article}
\usepackage[T1]{fontenc}
\usepackage[utf8]{inputenc}
\usepackage[french]{babel}
\usepackage{amsmath,amssymb,mathrsfs,mathpazo}
\usepackage{geometry}
\geometry{a4paper,margin=2cm}
\usepackage{graphicx,array,tabularx,tabularray,colortbl,multirow,multicol,booktabs,enumitem,pifont,ulem,fancybox,diagbox,cancel}
\usepackage{pstricks,pstricks-add,pst-plot,pst-tree,pst-node,pst-text,pst-func,pst-math,pst-bspline,pst-3dplot,pst-eucl,pst-solides3d}
\let\pscalculate\undefined
\usepackage{pst-bezier}
\usepackage{tikz,pgfplots,tkz-tab}
\usetikzlibrary{arrows,arrows.meta,patterns,calc,positioning,trees,decorations.markings,shapes.geometric,shapes.misc,shapes.arrows,intersections,decorations.pathmorphing,decorations.pathreplacing,backgrounds,shadows,babel}
\pgfplotsset{compat=1.18}
\usepgfplotslibrary{fillbetween}
\usepackage[np]{numprint}
\usepackage{listings,makecell,pas-tableur,twoopt}
\let\barre\undefined
\usepackage[pstricks]{tablvar}
\definecolor{darkgreen}{RGB}{0,100,0}
\providecommand{\red}{\color{red}}
\providecommand{\blue}{\color{blue}}
\providecommand{\green}{\color{green}}
\providecommand{\black}{\color{black}}
\providecommand{\white}{\color{white}}
\providecommand{\gray}{\color{gray}}
\providecommand{\cyan}{\color{cyan}}
\providecommand{\magenta}{\color{magenta}}
\providecommand{\yellow}{\color{yellow}}
\providecommand{\violet}{\color{violet}}
\providecommand{\orange}{\color{orange}}
\providecommand{\euro}{\texteuro}
\providecommand{\EUR}{\texteuro}
\providecommand{\ieme}{\textsuperscript{e}}
\providecommand{\degres}{\ensuremath{^\circ}}
\providecommand{\degree}{\ensuremath{^\circ}}
\providecommand{\vv}[1]{\overrightarrow{#1}}
\providecommand{\Coord}[1]{#1}
\providecommand{\decosix}{}
\providecommand{\starredbullet}{\ensuremath{\bullet}}
\pagestyle{empty}
\setlength{\parindent}{0pt}
\usepackage{adjustbox}\DeclareUnicodeCharacter{2236}{:}\DeclareUnicodeCharacter{2190}{\ensuremath{\leftarrow}}
\providecommand{\R}{\mathbb{R}}
\providecommand{\N}{\mathbb{N}}
\providecommand{\Z}{\mathbb{Z}}
\providecommand{\Q}{\mathbb{Q}}
\providecommand{\D}{\mathbb{D}}
\providecommand{\C}{\mathbb{C}}
\providecommand{\e}{{\mathrm{e}}}
\providecommand{\vect}[1]{\overrightarrow{#1}}
\providecommand{\vectt}[1]{\overrightarrow{\mathrm{#1}}}
\providecommand{\vvt}[1]{\overrightarrow{\mathrm{#1}}}
\providecommand{\np}[2][]{{#2~#1}}
\providecommand{\vv}[1]{{\overrightarrow{#1}}}
\renewcommand{\d}{{\mathrm{d}}}
\providecommand{\degres}{{^\circ}}
\providecommand{\degre}{{^\circ}}
\providecommand{\blue}{\color{blue}}
\providecommand{\red}{\color{red}}
\providecommand{\enskip}{\quad}
\providecommand{\llbracket}{\lbrack\!\lbrack}
\providecommand{\rrbracket}{\rbrack\!\rbrack}
\providecommand{\euro}{{\text{€}}}
\providecommand{\no}{{\text{n°}}}
\providecommand{\interval}[2]{{[#1;#2]}}
\providecommand{\intervFO}[2]{{[#1;#2[}}
\providecommand{\dotfill}{\ldots}

\providecommand{\ds}{\displaystyle}
\providecommand{\textding}[1]{#1}
\providecommand{\Oij}{\(\mathrm{O};\vec{\imath},\vec{\jmath}\)}
\providecommand{\Oijk}{\(\mathrm{O};\vec{\imath},\vec{\jmath},\vec{k}\)}
\providecommand{\Ouv}{\(\mathrm{O};\vec{u},\vec{v}\)}
\providecommand{\intervFF}[2]{\ensuremath{[#1;#2]}}
\renewcommand{\euro}{\ifmmode\text{\texteuro}\else\texteuro\fi}
\DeclareUnicodeCharacter{00B0}{\ifmmode^{\circ}\else\textdegree\fi}
\DeclareUnicodeCharacter{20AC}{\ifmmode\text{\texteuro}\else\texteuro\fi}
% Énoncés : Éducation nationale. Transcription et corrigés : APMEP, Denis Vergès et contributeurs. Découpage et classement : CoopMaths, Sébastien Lozano, Éric Elter et contributeurs. Adaptation de la présentation : Maths Explorer.
% CC BY-SA : https://coopmaths.fr/www/about/
% https://coopmaths.fr/alea/static/eam/2026/tex/eam_technologique_2026_06_etrangers_automatismes.tex
\begin{document}
\noindent\textit{APMEP / CoopMaths. CC BY-SA. Présentation adaptée par Maths Explorer.}\par\medskip
\textbf{Pour cette première partie, aucune justification n'est demandée et une seule réponse est possible par question. Pour chaque question, reportez son numéro sur votre copie et indiquez votre réponse.}

Une réponse fausse ou l'absence de réponse n'enlève aucun point.

\subsubsection*{Question 1}
$5-\dfrac{3}{2}$ est égal à :

\textbf{a.} $1$ \hfill \textbf{b.} $\dfrac{7}{2}$ \hfill \textbf{c.} $4$ \hfill \textbf{d.} $-\dfrac{5}{2}$

\subsubsection*{Question 2}
$\left(2^3\right)^{-5}$ est égal à :

\textbf{a.} $2^{-2}$ \hfill \textbf{b.} $2^{-15}$ \hfill \textbf{c.} $6^{-5}$ \hfill \textbf{d.} $2^8$

\subsubsection*{Question 3}
Un téléphone coûte 990 euros. Le prix baisse de 20 \%.

Son nouveau prix est :

\textbf{a.} $990\times 0{,}2$ \hfill \textbf{b.} $990\times\left(1+\dfrac{20}{100}\right)$ \hfill \textbf{c.} $990\times\left(-\dfrac{20}{100}\right)$ \hfill \textbf{d.} $990\times 0{,}8$

\subsubsection*{Question 4}
On considère un dé truqué tel que la probabilité d'obtenir un 5 et celle d'obtenir un 6 sont chacune égales à 0,3.

On lance le dé. La probabilité d'obtenir un nombre supérieur ou égal à 5 est :

\textbf{a.} $\dfrac{2}{6}$ \hfill \textbf{b.} $\dfrac{1}{6}$ \hfill \textbf{c.} $0{,}6$ \hfill \textbf{d.} $0{,}3$

\subsubsection*{Question 5}
Les solutions de l'équation $(x-2)(2x+1)=0$ sont :

\textbf{a.} $2$ et $-\dfrac{1}{2}$ \hfill \textbf{b.} $-2$ et $1$ \hfill \textbf{c.} $2$ et $\dfrac{1}{2}$ \hfill \textbf{d.} $2$ et $-1$

\subsubsection*{Question 6}
Voici les notes obtenues dans une classe lors d'un contrôle en mathématiques.
\begin{center}
\begin{tabular}{|c|c|c|c|c|}
\hline
Note & 7 & 10 & 12 & 14\\
\hline
Nombre d'élèves & 5 & 7 & 8 & 10\\
\hline
\end{tabular}
\end{center}
La note médiane de ce contrôle est :

\textbf{a.} $12$ \hfill \textbf{b.} $11$ \hfill \textbf{c.} $11{,}37$ \hfill \textbf{d.} $7{,}5$

\subsubsection*{Question 7}
On donne ci-dessous la courbe représentative $\mathcal{C}_f$ d'une fonction $f$ définie sur l'intervalle $[-4;3]$.
\begin{center}
\begin{tikzpicture}[scale=0.78]
\draw[step=0.5,gray!35,thin] (-4.5,-1.5) grid (3.5,5.5);
\draw[->,thick] (-4.5,0)--(3.6,0);
\draw[->,thick] (0,-1.55)--(0,5.6);

\foreach \x in {-4,-3,-2,-1,0,1,2,3} {\draw (\x,0.07)--(\x,-0.07) node[below,scale=0.75] {\x};}
\foreach \y in {-1,1,2,3,4,5} {\draw (0.07,\y)--(-0.07,\y) node[left,scale=0.75] {\y};}

\begin{scope}
\clip (-4.5,-1.5) rectangle (3.5,5.5);

\draw[gray!75,very thick,domain=-4:3,samples=200,smooth] 
    plot (\x, {-1/6*pow(\x,3) + 13/6*\x + 2});
\end{scope}

\node at (2.75,5.15) {$\mathcal{C}_f$};


\fill (-4,4) circle (1.5pt);  
\fill (-3,0) circle (1.5pt);  
\fill (-2,-1) circle (1.5pt);  
\fill (-1,0) circle (1.5pt);  
\fill (0,2) circle (1.5pt);   
\fill (2,5) circle (1.5pt);   
\fill (3,4) circle (1.5pt);   
\end{tikzpicture}
\end{center}
L'ensemble des solutions de l'inéquation $f(x)\geqslant 0$ est :

\textbf{a.} $[0\,;\,3]$ \hfill \textbf{b.} $[-4\,;\,-2]$ \hfill \textbf{c.} $[-4\,;\,-3]\cup[-1\,;\,3]$ \hfill \textbf{d.} $[-4\,;\,-3]\cup[0\,;\,3]$

\subsubsection*{Question 8}
Dans le lycée Alpha, il y a 500 élèves.

150 lycéens pratiquent un sport.

Le pourcentage d'élèves pratiquant un sport dans ce lycée est égal à :

\textbf{a.} $15\ \%$ \hfill \textbf{b.} $35\ \%$ \hfill \textbf{c.} $30\ \%$ \hfill \textbf{d.} $65\ \%$




\begin{center}
\end{center}

\end{document}
